% ECONOMICS_PROBLEM: {"id": "C90-646", "title": "Experimenter-demand and Hawthorne effects", "jel_code": "C90", "source_id": "OP-0068"}
% !TeX program = xelatex
\documentclass[11pt,a4paper]{article}
\usepackage[margin=23mm]{geometry}
\usepackage{fontspec}
\setmainfont{Latin Modern Roman}
\usepackage{amsmath,amssymb,mathtools}
\usepackage{xcolor,enumitem,booktabs,longtable,array,xurl,hyperref}
\definecolor{muted}{HTML}{53636C}
\hypersetup{colorlinks=true,urlcolor=blue,linkcolor=blue}
\setlength{\parindent}{0pt}
\setlength{\parskip}{5pt}
\newcommand{\R}{\mathbb R}
\newcommand{\E}{\mathbb E}
\newcommand{\N}{\mathbb N}
\newcommand{\ind}{\mathbf 1}
\newcommand{\Var}{\operatorname{Var}}
\newcommand{\Cov}{\operatorname{Cov}}
\newcommand{\supp}{\operatorname{supp}}
\DeclareMathOperator*{\argmin}{arg\,min}
\DeclareMathOperator*{\argmax}{arg\,max}
\newcommand{\entrymeta}[3]{{\sffamily\small\color{muted}\textbf{\texttt{#1}}\quad|\quad #2\par #3\par}}
\newcommand{\Problemref}[1]{\texttt{#1}}
\newcommand{\JELref}[2]{\texttt{#2} (JEL \texttt{#1})}
\begin{document}
\section*{Experimenter-demand and Hawthorne effects}
\textbf{Problem ID:} \texttt{C90-646}\quad \textbf{JEL:} \texttt{C90}\par
% BEGIN ECONOMICS STATEMENT
\label{problem:OP-0068}
\label{atlas:prob:B8}

\noindent\textbf{Original identifier:} B8 (atlas item 68). \textbf{Classification:} Experimental measurement and partial identification research program. \textbf{Original tractability label:} Easy (editorial, not a complexity result).

\subsection*{Definitions and assumptions}

Let $A\in\{0,1\}$ be treatment and $D\in\{-,0,+\}$ a randomized message intended to change perceived researcher expectations. Let $H\in\{0,1\}$ indicate awareness of being observed. Define $Y(a,d,h)$ with consistency, no interference, factorial randomization, and a bounded or integrable outcome scale. An effect of researcher demand is a response to perceived expectations; an observation-awareness effect can occur without any inferred desired action. The administrative outcome used for verification must have comparable measurement across arms. The neutral message $D=0$ need not remove perceived demand.

\subsection*{Formal research question}

Identify $\tau(d,h)=\mathbb E[Y(1,d,h)-Y(0,d,h)]$ and demand-message contrasts. Determine assumptions sufficient to bound the target effect without experimental demand,
\[
\tau^*=\mathbb E[Y^*(1)-Y^*(0)],
\]
using randomized $D$ arms. For example, if for each $a$, messages bracket demand-free mean outcomes as
\[
\mathbb E[Y(a,-,1)]\leq\mathbb E[Y^*(a)]\leq\mathbb E[Y(a,+,1)],
\]
then derive the implied lower and upper bounds for $\tau^*$ and valid joint uncertainty regions. Establish whether the bracket assumptions are credible under independently measured beliefs about researcher intent.

\subsection*{Interpretation and scope}

Factorial randomization identifies effects of the assigned messages, not automatically an abstract psychological demand effect. A message may supply factual information or change trust, violating the interpretation that it changes only perceived expectations. Conversely, a null message contrast can arise because subjects do not believe the manipulation or because demand operates similarly in all arms. The bracketing restriction is additional substantive content; it is not a consequence of randomized assignment. Demand-free outcomes may remain counterfactual even when administrative behavior is observed unobtrusively, because initial recruitment can alter behavior. Studies should specify whether they seek an awareness-free, message-free, or ordinary-policy effect. Repeated manipulation may itself increase scrutiny, so between-person designs can be preferable. Generalizing bounds across tasks requires a transportability assumption and outcome-scale normalization. Routine inclusion of a placebo arm is useful only when its exclusion restrictions and inferential target are made explicit.

\subsection*{Source audit and status}

[S1] reviews economic demand effects; [S2] develops measurement and bounds; [S3] finds limited demand responses in particular survey experiments. The 2019 issue year of [S3] is compatible with its 2018 online publication. These verified sources do not support a universal assertion that experimental results are mostly demand-driven. The distinction between demand and observation awareness corrects the original conflation; adapting the bracketing assumptions to specified economic settings is an editorial extension.

\subsection*{References}

\begin{enumerate}[label={[S\arabic*]},leftmargin=*]

\item Daniel John Zizzo (2010). \emph{Experimenter Demand Effects in Economic Experiments}. Experimental Economics 13(1), 75--98. \url{https://doi.org/10.1007/s10683-009-9230-z}\par \textit{Source role:} Methodological review; article year differs from platform migration date. \textit{Locator:} Abstract and article metadata.

\item Jonathan de Quidt, Johannes Haushofer, and Christopher Roth (2018). \emph{Measuring and Bounding Experimenter Demand}. American Economic Review 108(11), 3266--3302. \url{https://www.aeaweb.org/articles?id=10.1257/aer.20171330}\par \textit{Source role:} Demand manipulation and bounding method. \textit{Locator:} Abstract and article metadata.

\item Jonathan Mummolo and Erik Peterson (2019). \emph{Demand Effects in Survey Experiments: An Empirical Assessment}. American Political Science Review 113(2), 517--529. \url{https://doi.org/10.1017/S0003055418000837}\par \textit{Source role:} Evidence on specific survey experiments. \textit{Locator:} Abstract and article metadata.

\end{enumerate}

\subsection*{Common conventions: Source mathematical conventions}

Unless an entry states otherwise, agent and firm sets are finite. State and action spaces are measurable, strategies and outcome maps are measurable, probability laws are specified, and expectations exist for the targets under discussion. Infinite-horizon discount factors lie in $(0,1)$ and discounted objectives are finite. Norms, tolerances and complexity input sizes must be specified for computational comparisons. Constants or primitives that appear locally are inputs, not empirical facts asserted by this catalogue.

\textbf{Identification.} A statistical model $m\in\mathcal M$ implies an observable law $P_m$ and target $\theta(m)$. At law $P$, its identified set is
\[
\Theta(P;\mathcal M)=\{\theta(m):m\in\mathcal M,\ P_m=P\}.
\]
Point identification holds when this set is a singleton, or equivalently when $P_{m_1}=P_{m_2}$ implies $\theta(m_1)=\theta(m_2)$ for all compatible models. Sharp bounds mean bounds attained, or approached when the boundary is not attained, by compatible models. Writing this set is a definition, not a solution: the substantive task is to establish informative restrictions and derive or estimate the set. An unrestricted-confounding model may give only trivial bounds.

\textbf{Inference.} Identification, estimability and valid uncertainty quantification are separate questions. Fix $\alpha\in(0,1)$. A confidence procedure $C_n$ over a declared law class $\mathcal P$ has asymptotic uniform coverage at level $1-\alpha$ when
\[
\liminf_{n\to\infty}\inf_{P\in\mathcal P}\Pr_P\{\theta(P)\in C_n\}\geq1-\alpha.
\]
For set-identified targets the coverage event must be defined separately, for example coverage of the full identified set. If an entry uses identification-indexed models instead of uniquely indexed laws, the same coverage convention is applied over those models. Pointwise limits do not establish uniform validity, and asymptotic validity does not guarantee finite-sample precision. Sample sizes, dependence structures and weak-identification sequences are part of each application.

\textbf{Counterfactuals.} $Y_i(a)$ is a potential outcome under a specified intervention, including its timing, implementation and versions. It is not $Y_i$ conditional on voluntarily choosing $a$. A full intervention vector permits interference; shorthand scalar treatment notation does not silently impose no interference. Assignment, selection, attrition, anticipation and measurement must be specified. A claim about an instrument's local effect is not automatically a population policy effect.

\textbf{Economic equilibrium.} A counterfactual model includes best responses, constraints, feasibility, market clearing, state transitions and expectations consistent with the stated information. When several equilibria exist, either a declared selector is an input or the target is set-valued. Comparisons across policy regimes must use the stated selection convention. Reduced-form relationships are not presumed invariant after an equilibrium-changing intervention.

\textbf{Optimization and welfare.} Optimization is over the stated feasible set. If attainment is not established, $\sup$ or $\inf$ and $\varepsilon$-optimal choices replace an asserted maximizer or minimizer. For example, with value $V=\sup_{a\in\mathcal A}W(a)<\infty$, an $\varepsilon$-optimal set is $\{a\in\mathcal A:W(a)\geq V-\varepsilon\}$ for $\varepsilon>0$. Benefits, costs, risk and health or environmental outcomes can be aggregated only using declared units and conversion weights. Interpersonal utility weights, the treatment of fiscal transfers and foreign profits, discounting, and the behavioral welfare criterion are explicit inputs. A comparative-static derivative additionally requires existence and appropriate differentiability of the selected counterfactual.

\textbf{Sources.} Local labels [S1], [S2], and so forth restart in every question. Locators identify the relevant abstract, model, section or result at the level actually checked; a bibliographic or abstract-level verification is not represented as inspection of an entire proof. Working papers are distinguished from later publications and changing versions. Source summaries are paraphrased. All URL checks and editorial formulations refer to the audit date stated above.
% END ECONOMICS STATEMENT
\end{document}
